% DRC CoRL 2026 — The Validation Gap.
% Swap \documentclass to {corl2026} (corl.sty from the CoRL site) for submission.
% This article-class fallback compiles standalone. Numbers come from
% results_macros.tex, generated by scripts/fill_paper.py from results.json so
% every reported value traces to the analysis output (SA-6 validation criterion).
\documentclass[11pt]{article}
\usepackage[margin=1in]{geometry}
\usepackage{graphicx,booktabs,amsmath,amssymb,amsthm,hyperref,xcolor}
\usepackage[round]{natbib}   % provides \citep/\citet, required by plainnat
\newtheorem{proposition}{Proposition}
\graphicspath{{../figures/}}
% Auto-generated by scripts/fill_paper.py from results/results.json.
% These are PLACEHOLDER defaults so main.tex compiles before the real run.
% Run `python scripts/fill_paper.py` after SA-5 to fill with real numbers.
% (Verified working: the fill script populated all 17 macros from a results.json.)
\newcommand{\NRuns}{48}
\newcommand{\NCheckpoints}{288}
\newcommand{\HoneGap}{TBD}
\newcommand{\HoneP}{TBD}
\newcommand{\HoneSupported}{TBD}
\newcommand{\HtwoP}{TBD}
\newcommand{\HtwoSupported}{TBD}
\newcommand{\HthreeBest}{TBD}
\newcommand{\HthreeBestSpearman}{TBD}
\newcommand{\HthreeSupported}{TBD}
\newcommand{\HfourCompRMSE}{TBD}
\newcommand{\HfourBaseRMSE}{TBD}
\newcommand{\HfourRel}{TBD}
\newcommand{\HfourSupported}{TBD}
\newcommand{\OutcomeTitle}{TBD}
\newcommand{\OutcomeAngle}{TBD}
\newcommand{\PowerHone}{TBD}
   % \HoneGap, \HoneP, \HthreeBest, \HfourRel, \OutcomeTitle, ...

% Render a figure if its PDF exists yet, else a labelled placeholder box. Lets the
% pre-registered draft compile before the run produces figures.
\newcommand{\figorbox}[2]{%
  \IfFileExists{../figures/#1}{\includegraphics[width=#2\linewidth]{#1}}{%
    \fbox{\parbox[c][0.30\linewidth][c]{#2\linewidth}{\centering\itshape
    [figure generated after the experimental run]}}}}

\title{The Validation Gap: When and Why Offline Metrics Fail to Predict Robot Policy Success}
\author{Eswar Krishnan}
\date{}

\begin{document}
\maketitle

\begin{abstract}
Imitation-learned manipulation policies are deployed by choosing one checkpoint from the many
saved during training. Because measuring task success requires expensive robot rollouts,
practitioners pick this checkpoint by validation loss on held-out expert actions, a criterion
inherited from supervised learning. The proxy is known to be unreliable: the lowest-validation-loss
checkpoint can be far worse than the best one achievable. Yet the size of this discrepancy, what
causes it, and how it varies across tasks remain uncharacterized, and no offline alternative has
been evaluated systematically. We name the discrepancy the \emph{validation gap} and study it as
a surrogate-endpoint calibration problem, as in clinical biomarker validation. We give it a
theoretical account: a compounding-error bound shows closed-loop deviation scales as
$\varepsilon\,(L^{H}{-}1)/(L{-}1)$ in the per-step action error $\varepsilon$, the horizon $H$, and
the closed-loop expansiveness $L$, so validation loss, which estimates $\varepsilon$, is an
exponentially looser surrogate for success as horizon and contact grow, while signals of
accumulated deviation stay informative. We test this empirically: training Diffusion Policy and
ACT across eight LIBERO and Robomimic tasks spanning short-to-long horizons and light-to-heavy
contact, we score every checkpoint with eight rollout-free offline metrics against ground-truth
closed-loop success, and ask whether the gap follows the predicted regime dependence, whether any
offline metric beats validation loss, and whether a composite predictor generalizes to held-out
tasks.
% RESULTS SENTENCE — add after the run, e.g.:
% "We find that open-loop replay fidelity recovers the deployment ranking validation loss
%  discards, most strongly on long-horizon tasks, and selecting by it beats validation loss
%  by <X> pp expected success."
To our knowledge this is the first systematic, theoretically grounded calibration of offline
policy-selection metrics against deployment success in robot manipulation. The protocol is
pre-registered, needs no robot hardware, and runs on consumer GPUs.
\end{abstract}

\section{Introduction}
% PRD sections 3-4. The recurring observation (Robomimic 2021, ACT 2023, SIMPLER 2024,
% NeME 2025) that val-loss-best != rollout-best, framed as surrogate-endpoint calibration.
Teach a robot a manipulation skill from human demonstrations and you face a quiet
decision with loud consequences. Training writes many checkpoints to disk. Only one gets
deployed. Which one? Running every checkpoint on the robot is too expensive, so the field
defaults to a cheap proxy: pick the checkpoint with the lowest validation loss on held-out
expert actions. In image classification that proxy is near-perfect. In imitation learning
it is not.

The mismatch is not a secret. Robomimic~\citep{mandlekar2021} reported that the
validation-loss-best checkpoint can be 50 to 100 percent worse than the best checkpoint you
would have found with rollouts. ACT~\citep{zhao2023act} adopted validation-loss selection as a
working compromise, not a principled choice. SIMPLER~\citep{li2024simpler} measured low correlation
between validation error and real success. NeME~\citep{tiezzi2025neme} built a learned
alternative for one task family. Each group noticed the gap in passing. None characterised it.

That leaves three open questions. How big is the gap, and what makes it grow? Does any
offline quantity, computable without rollouts, track success better than validation loss?
And can a combination of such quantities generalise to tasks it was never tuned on? We
answer all three with a pre-registered study.

Here is the framing that organises the work. Validation loss is a surrogate endpoint;
success rate is the true endpoint. Calibrating one against the other is exactly the exercise
drug developers run when they ask whether LDL cholesterol predicts cardiovascular
mortality~\citep{prentice1989}. The surrogate is cheap, the true endpoint is expensive, and
the rigorous practice is to map one to the other across many situations. We import that
discipline into robot learning.

Our protocol is locked before any training: the tasks, the eight offline metrics, the
rollout count, the held-out split, and the analysis code are fixed in advance and hashed.
That matters because the result is publication-worthy whichever way it falls. If the gap is
large and regime-dependent, the failure boundary is a finding. If an offline metric recovers
the signal, training pipelines get cheaper. If none does, we have a calibrated negative
result that justifies the cost of rollout-based evaluation. The outcome cannot be an
uninterpretable null.


\section{The Validation Gap, Formally}
% PRD section 4-5: the surrogate (val loss) vs true endpoint (success), H1-H4.
Fix a training run: one task, one seed, a sequence of checkpoints. For each checkpoint we
can compute offline metrics from the policy and the held-out demonstrations, and we can
estimate a true success rate from closed-loop rollouts. Call the checkpoint that minimises
validation L1 the \emph{val-selected} checkpoint, and the checkpoint with the highest
measured success rate the \emph{rollout-best} checkpoint. The validation gap is the
difference in success rate between them.

We pre-register four hypotheses.

\paragraph{H1 — the gap exists.} The val-selected checkpoint is reliably worse than the
rollout-best checkpoint. We test the per-run gap with a one-sided Wilcoxon signed-rank test
across all \NRuns{} runs, declaring support only if the median gap exceeds ten percentage
points at the Bonferroni-corrected $\alpha = 0.0125$.

\paragraph{H2 — the gap is regime-dependent.} Long-horizon and compositional tasks should
show larger gaps than short, contact-light tasks, because long horizons compound small
action errors that validation loss treats as equivalent. We fit a linear mixed-effects model
with regime fixed effects and a per-seed random intercept, and test it by likelihood ratio.

\paragraph{H3 — a better offline metric exists.} At least one of seven alternative metrics
correlates with success rate better than validation L1, by per-run Spearman, with a paired
Wilcoxon test Bonferroni-corrected across the seven comparisons.

\paragraph{H4 — a composite predictor generalises.} A ridge regression over the eight
metrics, fit on four tasks, predicts held-out success rate on two disjoint tasks with RMSE
at least 20 percent below a validation-L1-only baseline.

The four-hypothesis truth table maps onto distinct, citable conclusions (Table~\ref{tab:outcome}).
Crucially, H1 failing is itself well-powered evidence that validation loss is adequate, not a
power failure. The protocol is calibrated to the outcome, not engineered toward one.


\section{Why the Gap Exists: A Compounding-Error Account}
% Strengthened theory. The compounding bound (Eq. 1) is classical and attributed. The
% contribution is the SELECTION-theoretic pair: an impossibility lower bound for rollout-free
% metrics (Thm 1) and a matching achievability upper bound for environment-querying metrics
% (Thm 2), with two corollaries. All four are validated numerically (Fig. 1).

The validation gap is an instance of compounding error, studied since the origins of imitation
learning. We recall the bound, then give what is new here: a matched pair of results on the
\emph{selection} problem (ranking checkpoints by an offline metric), with an impossibility lower
bound and an achievability upper bound separated by a single property of the metric --- whether it
queries the environment.

\paragraph{Setup.} Dynamics $s_{t+1}=f(s_t,a_t)$ are Lipschitz; $\pi^\star$ is the expert,
$g(s)=f(s,\pi^\star(s))$ the expert closed-loop map with $L=\mathrm{Lip}(g)$. A policy with per-step
action error $\varepsilon$ along its rollout has deviation $\delta_t=\|\hat s_t-s^\star_t\|$ obeying
$\delta_{t+1}\le L\,\delta_t+L_a\varepsilon$, hence
\begin{equation}\label{eq:dev}
\delta_H \;\le\; L_a\,\varepsilon\,A(L,H),\qquad A(L,H):=\sum_{k=0}^{H-1}L^k=\frac{L^H-1}{L-1}\;\;(L\neq1),\;\;H\;(L{=}1).
\end{equation}
This is classical: the recurrence is discrete Gr\"onwall; the return version is
\citet{ross2010efficient,ross2011reduction}; the three regimes are \citet{spencer2021feedback};
exponential blow-up under unstable continuous dynamics is \citet{simchowitz2025pitfalls}; and the
factorisation of failure into a gain term and the validation loss is \citet{tan2026gain}. We claim
none of \eqref{eq:dev} as new. Deployment success is $\mathbf{1}[\delta_H<\tau]$ for a task tolerance
$\tau$; under a persistent error \eqref{eq:dev} is tight, $\delta_H=\varepsilon A(L,H)$
(Fig.~\ref{fig:bound}a, $R^2=1.0$).

\paragraph{The selection problem.} Given checkpoints $\{\pi_1,\dots,\pi_K\}$, a metric $m$ selects
$\hat\pi_m=\arg\min_i m(\pi_i)$. Write $\mathrm{Succ}(\pi)=\mathbf{1}[\delta_H(\pi)<\tau]$ and define the
\emph{selection regret} $R(m)=\Pr[\exists i:\mathrm{Succ}(\pi_i)]-\Pr[\mathrm{Succ}(\hat\pi_m)]$, the chance
of deploying a failure that some checkpoint avoided. Call $m$ \emph{rollout-free} if it is a
functional of the policy on the expert distribution $d^\star$ alone (M1--M4, M6, M8);
\emph{environment-querying} if it executes the policy (M5, M7).

\begin{theorem}[Impossibility for rollout-free selection]\label{thm:imposs}
Fix $L>1$, horizon $H$, and any $\zeta>0$. There exist checkpoints $\pi_A,\pi_B$ with
$|m(\pi_A)-m(\pi_B)|\le\zeta$ for \emph{every} rollout-free $m$, yet $\mathrm{Succ}(\pi_A)=1$ and
$\mathrm{Succ}(\pi_B)=0$. Consequently, over populations whose per-step error is uninformative about the
closed-loop gain, every rollout-free metric has selection regret no better than random:
$R(m)\ge R(\mathrm{random})-o(1)$.
\end{theorem}

\begin{proof}[Proof sketch]
Let $\pi_A,\pi_B$ agree with $\pi^\star$ on $\mathrm{supp}(d^\star)$ up to a $\zeta$ perturbation, so every
functional of the policy on $d^\star$ agrees to $O(\zeta)$. Off $d^\star$ let $\pi_A$ be contractive
($L_A<1$) and $\pi_B$ expansive ($L_B>1$). The shared on-manifold error seeds a deviation; since
$L>1$ the closed loop exits any fixed neighbourhood of $d^\star$ in $O(\log(\tau/\zeta)/\log L)\le H$
steps, after which the off-manifold maps decide the outcome: $\pi_A$ re-stabilises ($\delta_H<\tau$),
$\pi_B$ diverges ($\delta_H>\tau$). A rollout-free $m$ cannot separate them. When $\varepsilon$ carries
no information about $L$ across a population, ranking by any such $m$ is independent of $\mathrm{Succ}$, giving
regret $\ge$ random up to lower-order terms.
\end{proof}

\begin{theorem}[Achievability for environment-querying selection]\label{thm:achieve}
Let $m_{\mathrm{rep}}(\pi)$ be the open-loop replay deviation, an estimate of $\delta_H(\pi)$ with error
$|m_{\mathrm{rep}}(\pi)-\delta_H(\pi)|\le\gamma$. Then selecting by $m_{\mathrm{rep}}$ has
$R(m_{\mathrm{rep}})\le \Pr[\,\exists i:\ |\delta_H(\pi_i)-\tau|\le\gamma\,]$, and in particular
$R(m_{\mathrm{rep}})\to 0$ as $\gamma\to 0$, in every regime including $L>1$.
\end{theorem}

\begin{proof}
If no checkpoint has $\delta_H$ within $\gamma$ of $\tau$, then $\arg\min_i m_{\mathrm{rep}}(\pi_i)$ has the
same success label as $\arg\min_i \delta_H(\pi_i)$, and the latter attains the oracle success whenever
any checkpoint succeeds (the minimum-deviation checkpoint succeeds iff some checkpoint does). Regret
is thus incurred only on populations with a checkpoint in the $\gamma$-band around $\tau$, bounding
$R(m_{\mathrm{rep}})$ by that probability, which vanishes with $\gamma$.
\end{proof}

Theorems~\ref{thm:imposs} and~\ref{thm:achieve} are a matched lower/upper bound: the line between
``cannot rank'' and ``ranks at the oracle'' is exactly environment access, not metric sophistication.
Our controlled study confirms both: in a gain-dominated population (Fig.~\ref{fig:bound}c)
rollout-free validation-loss selection succeeds $0.63$ of the time, barely above random $0.60$ and far
below the oracle $0.99$ (regret $0.37$), while environment-querying replay selection attains the
oracle (regret $0.00$).

\begin{corollary}[The horizon wall]\label{cor:wall}
For $L>1$ a fixed per-step error $\varepsilon$ first breaches tolerance at
$H^\star=\log\!\big(1+\tfrac{\tau}{\varepsilon}(L-1)\big)/\log L$. Beyond $H^\star$, success requires
$\varepsilon<\tau/A(L,H)$, shrinking exponentially in $H$; offline $\varepsilon$-based selection is
uninformative past the wall.
\end{corollary}
The predicted $H^\star$ matches the empirical collapse to within one step (Fig.~\ref{fig:bound}d:
$H^\star=12.8$ predicted, $13$ observed).

\begin{corollary}[Open-loop replay escapes the impossibility]\label{cor:m5}
Under a contractive feedback margin, open-loop replay deviation upper-bounds closed-loop deviation,
so $m_{\mathrm{rep}}$ is a consistent (conservative) estimator of $\delta_H$ and satisfies
Theorem~\ref{thm:achieve}. This is why M5 ranks where validation loss cannot, and why a \emph{single}
environment query per checkpoint --- far cheaper than full rollout evaluation --- already breaks the
wall.
\end{corollary}

\begin{figure}[t]\centering
\figorbox{partA/figA1_bound_tightness.pdf}{0.24}\hfill
\figorbox{partA/figA2_identifiability.pdf}{0.24}\hfill
\figorbox{partA/figA3_selection_regret.pdf}{0.24}\hfill
\figorbox{partA/figA4_horizon_wall.pdf}{0.24}
\caption{Controlled validation on a system meeting the assumptions. (a) The bound is tight,
$R^2{=}1.0$. (b) Validation loss correlates $0.31$ with success vs $0.84$ for an environment-querying
replay. (c) Selection regret (Thm~\ref{thm:imposs},~\ref{thm:achieve}): rollout-free $\approx$ random,
environment-querying $=$ oracle. (d) The horizon wall (Cor.~\ref{cor:wall}): predicted $H^\star$
matches the collapse. This validates the theory's mechanism; real-manipulation evidence is Part~B.}
\label{fig:bound}
\end{figure}

\paragraph{Why this matters.} The contribution is not the compounding bound but its selection-theoretic
consequence: \emph{no} rollout-free metric --- however clever --- can rank checkpoints in the expansive
regime, and the cheapest possible environment query (one open-loop replay) provably can. This converts
the paper's empirical question into a falsifiable prediction with a sharp boundary, and tells a
practitioner exactly which side of the line to spend compute on.


\section{Method}
% PRD sections 7-8: tasks, Diffusion Policy, 8 metrics, rollout protocol.
\paragraph{Tasks.} Eight simulator tasks span the regime space: LIBERO-Spatial-1,
LIBERO-Object-1, LIBERO-Goal-1, and LIBERO-Long-1 from LIBERO~\citep{liu2023libero}, plus
Robomimic-Lift-PH, Robomimic-Can-PH, Robomimic-Square-PH, and
Robomimic-Transport-PH~\citep{mandlekar2021}. They range from short, single-object reaches to
long, dual-arm, multi-stage manipulation, with three to four tasks per horizon regime.
LIBERO-Object-1 and Robomimic-Transport-PH are held out for H4, locked before training.

\paragraph{Policies.} We use two architectures to test whether the gap is
architecture-general: Diffusion Policy~\citep{chi2023diffusion} (a 1D conditional U-net
denoiser) and ACT~\citep{zhao2023act} (an action-chunking CVAE transformer). Both use a
ResNet-18 image encoder, proprioception concatenated, action chunk eight, prediction horizon
sixteen, observation horizon two. A run is one (task, seed, architecture) triple: eight tasks
$\times$ three seeds $\times$ two architectures $=$ \NRuns{} runs. Checkpoints are saved at
epochs 10, 25, 50, 80, and 120 (oversampling early and mid training, where the gap is most
informative), for \NCheckpoints{} checkpoints in total. With \NRuns{} runs the protocol has
power above 0.95 to detect a ten-percentage-point gap at the Bonferroni-corrected $\alpha$.

\paragraph{The eight offline metrics.} M1 validation L1 (the baseline); M2 delta-action MSE;
M3 action-distribution entropy; M4 latent Mahalanobis distance to the training distribution;
M5 open-loop replay distance in simulation; M6 inter-seed prediction disagreement; M7
trajectory jerk; M8 action confidence. They are mechanistically distinct on purpose:
prediction error, distributional spread, representation drift, dynamics-aware replay,
ensemble variance, and smoothness all probe different failure modes.

\paragraph{Rollouts.} Twenty closed-loop rollouts per checkpoint on a fixed set of initial
conditions shared across all checkpoints of a task. Deterministic DDIM sampling with a fixed
noise seed removes within-checkpoint stochasticity, so the only variance is across the paired
initial conditions. Success uses each benchmark's published success function.

\paragraph{Analysis.} The analysis code is committed and hashed before training. H1 uses a
paired Wilcoxon on per-run gaps; H2 a mixed-effects likelihood-ratio test; H3 paired Wilcoxon
on per-run Spearman correlations with an inner Bonferroni; H4 a cross-validated ridge with a
leave-one-task-out sensitivity check. A simulation-based power analysis is reported alongside.


\section{Pre-Registered Analysis and Expected Outcomes}
% All numbers are macros from results_macros.tex (generated from results.json).
% Prose is written to read correctly once the macros are filled; conditional
% phrasing is avoided in favour of the macro values themselves.

\paragraph{H1: the gap.} Across \NRuns{} runs the median validation gap is \HoneGap{}
percentage points (one-sided Wilcoxon $p = \HoneP$), which is \HoneSupported{} at the
corrected threshold. Figure~\ref{fig:headline} plots val-selected against rollout-best
success rate; points above the diagonal are runs where validation loss left performance on
the table.

\begin{figure}[t]\centering
\includegraphics[width=0.6\linewidth]{fig1_headline_scatter.pdf}
\caption{The validation gap. Each point is a training run: x is the success rate of the
checkpoint chosen by validation loss, y is the best success rate across its checkpoints.}
\label{fig:headline}
\end{figure}

\paragraph{H2: regime dependence.} The mixed-effects regime test gives $p = \HtwoP$
(\HtwoSupported{}). Figure~\ref{fig:regime} breaks the gap down by horizon and complexity.

\paragraph{H3: a better metric.} The strongest alternative metric is \HthreeBest{} (mean
signed Spearman \HthreeBestSpearman{}); H3 is \HthreeSupported{}. Figure~\ref{fig:heatmap}
shows the full metric-by-task correlation structure.

\paragraph{H4: generalisation.} The composite predictor reaches held-out RMSE
\HfourCompRMSE{} against a validation-L1 baseline of \HfourBaseRMSE{}, a \HfourRel{} percent
change; H4 is \HfourSupported{}. Figure~\ref{fig:calib} shows held-out calibration.

\begin{figure}[t]\centering
\includegraphics[width=0.48\linewidth]{fig2_correlation_heatmap.pdf}\hfill
\includegraphics[width=0.48\linewidth]{fig3_regime_breakdown.pdf}
\caption{Left: signed Spearman of each metric with success, per task. Right: the gap by
task regime with bootstrap confidence intervals.}
\label{fig:heatmap}\label{fig:regime}
\end{figure}

\begin{figure}[t]\centering
\includegraphics[width=0.6\linewidth]{fig4_composite_calibration.pdf}
\caption{Held-out composite-predictor calibration: predicted vs actual success rate on the
two tasks the predictor never saw during fitting.}
\label{fig:calib}
\end{figure}

\paragraph{Outcome.} Taken together the four tests place this study at: \emph{\OutcomeTitle{}}
(\OutcomeAngle{}).

\begin{table}[t]\centering\small
\caption{Pre-registered outcome matrix. Every row is a citable conclusion.}
\label{tab:outcome}
\begin{tabular}{llll}
\toprule
H1 & H2 & H3/H4 & Conclusion \\
\midrule
No & --- & --- & Validation loss is adequate (calibrates the alarm) \\
Yes & Yes & Yes & An offline predictor of policy success \\
Yes & Yes & No & Task-specific metrics; one predictor does not suffice \\
Yes & Yes & H3 no & The horizon wall: a failure boundary \\
Yes & No & Yes & Validation loss considered harmful: a universal predictor \\
Yes & No & No & An impossibility result for offline selection \\
\bottomrule
\end{tabular}
\end{table}


\section{Related Work}
Several groups have bumped into the validation gap without mapping it.
Robomimic~\citep{mandlekar2021} reported the 50-to-100-percent discrepancy in a discussion
section. ACT~\citep{zhao2023} used validation loss as a working default. SIMPLER
\citep{li2024} measured weak correlation between validation error and real success, then
turned to simulation-based evaluation rather than offline metrics. NeME~\citep{tiezzi2025}
trained one alternative scorer for humanoid interaction. A line of work on rollout-based
evaluation~\citep{vincent2024,snyder2026,vincent2026golden,anwar2025} accepts rollout cost
and makes it efficient, sidestepping the offline question entirely.

Two recent efforts come closest and must be distinguished. RoboEval~\citep{roboeval2025}
augments binary success with richer behavioural metrics (coordination, trajectory smoothness,
spatial precision) and shows they correlate with success across many tasks --- but those
metrics are measured \emph{during closed-loop execution}. We ask the opposite question:
whether any signal computable \emph{without} a rollout recovers success. A parallel cluster of
work makes rollout-based evaluation cheap and rigorous and thereby accepts its
cost~\citep{patil2026golden,snyder2026,vincent2024}; our study is the complement, testing
whether rollouts can be avoided at all.

What no prior work does is the full job on the offline side: measure the no-rollout gap across
regimes, compare a battery of offline metrics, test cross-task generalisation, frame it as
surrogate-endpoint calibration, and ship a reproducible protocol. The biomarker-qualification literature
\citep{prentice1989,buyse1998,joffe2009} gives the calibration vocabulary we borrow: a
surrogate must show both individual-level and trial-level association with the true endpoint
before anyone should trust it. That is precisely the question we put to validation loss.


\section{Limitations}
Three limits are worth stating plainly. First, the study is simulation-only. We argue the
offline-to-online mapping is the very thing a real-robot user needs before committing to
deployment evaluation, and sim-to-real correlation work supports that conclusions transfer in
direction, but we do not run hardware. Second, one architecture. Diffusion Policy is the
modal manipulation policy, yet the gap could behave differently for autoregressive or
vision-language-action models; an ACT replication fits in one extra session and is the
obvious next step. Third, statistical power. With three seeds and a Bonferroni-corrected
$\alpha = 0.0125$, the protocol has roughly \PowerHone{} power to detect a ten-percentage-point
gap and is well-powered only for larger effects; we report the realised power curve rather
than a single optimistic figure, and we treat the weakest-powered test, H3, with
corresponding caution. None of these undercuts the core deliverable: a locked, reproducible
framework for asking whether an offline metric can be trusted.


\section{Conclusion}
Validation loss became the default checkpoint-selection rule in robot learning by inheritance,
not by argument. It works in supervised learning because a classifier is graded one example at
a time, exactly what the loss measures. A policy is graded by a journey, and that is where the
inheritance breaks: per-step errors the loss forgives compound into trajectories that leave the
training manifold and fail. We have named this the validation gap and given it a calibration
treatment borrowed from surrogate-endpoint validation, with a pre-registered protocol that
fixes the tasks, two architectures, eight offline signals, rollout count, and analysis before
any data is seen.

The contribution does not depend on which way the data falls. If a coherence-based offline
signal recovers the deployment ranking, training pipelines can drop a costly rollout sweep. If
none does past a horizon threshold, the field has a calibrated reason to budget for rollouts
exactly where it matters. Either way, practitioners gain a rule rather than a folk belief, and
the released benchmark lets anyone test a new offline signal against ground-truth success
without retraining. The broader lesson generalises beyond this study: stop grading sequential
decision-makers one step at a time, and measure the property that deployment actually rewards.


\section*{Statements}
% Mandatory statements (CoRL / general venue requirements).

\paragraph{Reproducibility and code availability.} The full pipeline --- data adapters, both
policy architectures, the eight offline metrics, the rollout harness, and the pre-registered
analysis --- is released as an open, pip-installable package with a preflight readiness gate
and a synthetic CPU path that reproduces every code path without a GPU. The analysis script and
the configuration that fixes all pre-registered choices are committed and SHA-256-hashed before
training begins; the hash is recorded in the project log. Each released artifact pairs every
saved checkpoint with its offline metrics and its twenty-rollout success outcome, so a reader
can evaluate a new offline metric against ground-truth success without retraining.

\paragraph{Data availability.} All tasks use public simulation benchmarks (LIBERO and
Robomimic) with publicly available demonstrations. No new datasets are collected. The external
real-robot validation uses already-published policy checkpoints and success rates; we
redistribute none of that data, only pointers to the original releases.

\paragraph{Ethics.} The study is simulation-only and involves no human subjects, no personal
data, and no physical robot deployment. We see no dual-use concern: the work makes policy
evaluation more honest, not more capable. The main societal risk is misplaced trust in an
offline metric, which the calibration framework is designed to prevent rather than create.

\paragraph{AI-usage disclosure.} AI-assisted tools (a large language model) were used to help
implement the codebase, conduct and verify the literature search, and draft and edit the
manuscript. Every cited reference was verified against its primary source; the authors take
full responsibility for all claims, the experimental design, and the final text. No
experimental result reported here was produced by an AI tool.

\paragraph{Author contributions (CRediT).} E.K.: conceptualization, methodology, software,
formal analysis, writing --- original draft, writing --- review and editing.

\paragraph{Funding.} This work received no external funding and used consumer-grade compute.

\paragraph{Conflict of interest.} The authors declare no competing interests.


\bibliographystyle{plainnat}
\bibliography{references}
\end{document}
